\documentclass{article}
\usepackage{amsmath,amssymb,amsthm}
\usepackage{algorithm,algorithmic}
\usepackage{fullpage}
\usepackage{hyperref}
\newtheorem{theorem}{Theorem}
\newtheorem{lemma}{Lemma}
\newtheorem{assumption}{Assumption}
\newcommand{\E}{\mathbb{E}}
\newcommand{\R}{\mathbb{R}}

\begin{document}

%======================================================================
\section*{Quantized Stochastic Gradient Descent (QSGD)}
%======================================================================

\section*{Mathematical Formulation}
Let $f:\R^{d}\rightarrow\R$ be the objective function and assume that we have access to a stochastic gradient oracle
\[
g_{t}= \nabla f(w_{t};\xi_{t}),
\]
where $\xi_{t}$ is a random data sample drawn i.i.d. from a distribution $\mathcal{D}$.  
A (possibly lossy) \emph{quantization operator}
\[
Q:\R^{d}\rightarrow\R^{d}
\]
satisfies the following two properties:

\begin{enumerate}
\item \textbf{Unbiasedness:} \(\displaystyle \E\!\big[\,Q(g)\mid g\big]=g.\)
\item \textbf{Bounded variance:} there exists a constant \(\omega\ge 0\) such that
\[
\E\!\big[\|Q(g)-g\|^{2}\mid g\big]\le \omega\|g\|^{2}\qquad\forall g\in\R^{d}.
\]
\end{enumerate}

The Quantized SGD iteration is
\[
\boxed{w_{t+1}=w_{t}-\eta_{t}\,Q\!\big(g_{t}\big)}\tag{1}
\]
with a (possibly time‑varying) stepsize \(\eta_{t}>0\).

\section*{Pseudocode}
\begin{algorithm}[H]
\caption{Quantized Stochastic Gradient Descent (QSGD)}
\begin{algorithmic}[1]
\STATE {\bf Input:} initial point $w_{0}\in\R^{d}$, stepsize schedule $\{\eta_{t}\}_{t\ge0}$, number of iterations $T$.
\FOR{$t=0$ {\bf to} $T-1$}
    \STATE Sample $\xi_{t}\sim\mathcal{D}$ and compute stochastic gradient $g_{t}= \nabla f(w_{t};\xi_{t})$.
    \STATE Quantize: $\tilde g_{t}= Q(g_{t})$.
    \STATE Update: $w_{t+1}= w_{t}-\eta_{t}\tilde g_{t}$.
\ENDFOR
\STATE {\bf Output:} $w_{T}$ (or the average $\bar w_{T}= \frac1T\sum_{t=0}^{T-1}w_{t}$).
\end{algorithmic}
\end{algorithm}

\section*{Dry Run Example}
Consider the scalar quadratic
\[
f(w)=(w-3)^{2},\qquad \nabla f(w)=2(w-3).
\]
We use a deterministic “quantizer’’ that simply rounds to the nearest multiple of $0.05$:
\(Q(g)=\operatorname{round}_{0.05}(g)\).  
Clearly \(\E[Q(g)]=g\) (deterministic) and \(\omega=0\).

\begin{center}
\begin{tabular}{c|c|c|c|c}
$t$ & $w_{t}$ & $g_{t}=2(w_{t}-3)$ & $Q(g_{t})$ & $w_{t+1}=w_{t}-0.1\,Q(g_{t})$\\ \hline
0 & $0.0$ & $-6.0$ & $-6.00$ & $0.0-0.1(-6)=0.6$\\
1 & $0.6$ & $-4.8$ & $-4.80$ & $0.6-0.1(-4.8)=1.08$\\
2 & $1.08$ & $-3.84$ & $-3.85$ & $1.08-0.1(-3.85)=1.465$\\
\end{tabular}
\end{center}

Corresponding objective values:
\[
f(0)=9,\; f(0.6)=5.76,\; f(1.08)=3.6864,\; f(1.465)\approx2.888.
\]
The iterates converge toward the minimizer $w^{*}=3$.

%======================================================================
\section*{Assumptions}
%======================================================================
\begin{assumption}[Convexity and Smoothness]\label{ass:convex}
The function $f:\R^{d}\rightarrow\R$ is convex and $L$‑smooth, i.e.
\[
f(y)\le f(x)+\langle\nabla f(x),y-x\rangle+\frac{L}{2}\|y-x\|^{2},\qquad\forall x,y\in\R^{d}.
\]
\end{assumption}

\begin{assumption}[Stochastic Gradient]\label{ass:sgd}
For all $t$,
\[
\E[g_{t}\mid w_{t}]=\nabla f(w_{t}),\qquad
\E\big[\|g_{t}-\nabla f(w_{t})\|^{2}\mid w_{t}\big]\le\sigma^{2},
\]
for some $\sigma^{2}\ge0$.
\end{assumption}

\begin{assumption}[Quantizer]\label{ass:quant}
The quantizer $Q$ satisfies unbiasedness and bounded relative variance with constant $\omega\ge0$:
\[
\E\!\big[Q(g)\mid g\big]=g,\qquad 
\E\!\big[\|Q(g)-g\|^{2}\mid g\big]\le\omega\|g\|^{2}.
\]
\end{assumption}

\begin{assumption}[Stepsize]\label{ass:stepsize}
The stepsize is chosen as
\[
\eta_{t}= \frac{c}{\sqrt{t+1}},\qquad c>0.
\]
\end{assumption}

%======================================================================
\section*{Main Convergence Theorem}
%======================================================================
\begin{theorem}[Expected Suboptimality of QSGD]\label{thm:main}
Under Assumptions \ref{ass:convex}–\ref{ass:stepsize}, let $w^{*}\in\arg\min_{w}f(w)$.  
Define the averaged iterate $\bar w_{T}= \frac{1}{T}\sum_{t=0}^{T-1}w_{t}$. Then
\[
\E\!\big[f(\bar w_{T})-f(w^{*})\big]\;\le\;
\frac{\|w_{0}-w^{*}\|^{2}}{2c\sqrt{T}}
\;+\;
\frac{cL}{2\sqrt{T}}\big(\sigma^{2}+ \omega G^{2}\big),
\tag{2}
\]
where $G^{2}\ge \sup_{t}\E\!\big[\|\nabla f(w_{t})\|^{2}\big]$ (finite for bounded domain).  
Consequently, $\E[f(\bar w_{T})-f(w^{*})]=\mathcal{O}(1/\sqrt{T})$.
\end{theorem}

\section*{Theoretical Properties}
\begin{itemize}
\item \textbf{Stability:} The recursion (1) remains stable for any $\eta_{t}\le 1/L$ because the smoothness inequality (Assumption~\ref{ass:convex}) bounds the growth of $\|w_{t+1}-w^{*}\|^{2}$.
\item \textbf{Effect of Quantization:} The extra term $\omega G^{2}$ in (2) quantifies the penalty incurred by the quantizer. When $\omega=0$ (exact communication) the bound reduces to the classical SGD rate.
\item \textbf{Hyper‑parameter Sensitivity:} The constant $c$ balances the bias–variance trade‑off: larger $c$ speeds early progress but inflates the variance term $\frac{cL}{2}\big(\sigma^{2}+\omega G^{2}\big)$.
\item \textbf{Comparison to Baselines:}
\begin{enumerate}
\item With exact gradients ($\omega=0$, $Q(g)=g$) we recover the standard $\mathcal{O}(1/\sqrt{T})$ SGD bound.
\item For deterministic gradients ($\sigma=0$) the rate becomes $\mathcal{O}(\omega/\sqrt{T})$, showing that quantization alone can dominate the error.
\end{enumerate}
\end{itemize}

\section*{Discussion}
The theorem shows that unbiased quantization does not alter the optimal $\mathcal{O}(1/\sqrt{T})$ convergence order; it merely inflates the constant by a factor $1+\omega G^{2}/\sigma^{2}$. This justifies the widespread use of low‑precision communication in distributed convex optimization.

%======================================================================
\section*{Preliminaries (Definitions and Lemmas)}
%======================================================================
\begin{lemma}[Smoothness Inequality]\label{lem:smooth}
If $f$ is $L$‑smooth then for any $x,y\in\R^{d}$,
\[
f(y)\le f(x)+\langle\nabla f(x),y-x\rangle+\frac{L}{2}\|y-x\|^{2}.
\]
\end{lemma}
\begin{proof}
Standard; follows from integrating the Lipschitz condition on $\nabla f$. \qedhere
\end{proof}

\begin{lemma}[Three‑Point Identity]\label{lem:three}
For any vectors $a,b,c\in\R^{d}$ and any scalar $\eta>0$,
\[
\|a-c\|^{2}= \|a-b\|^{2}+2\langle b-c,a-b\rangle+\|b-c\|^{2}.
\]
\end{lemma}
\begin{proof}
Expand $\|a-c\|^{2}=\|a-b+b-c\|^{2}$ and collect terms. \qedhere
\end{proof}

%======================================================================
\section*{Main Proof (Step‑by‑Step Derivation)}
%======================================================================
We prove Theorem~\ref{thm:main}.  Throughout the proof expectations are taken w.r.t. all randomness up to iteration $t$ unless otherwise noted.

\subsection*{Step 1: Distance Recursion}
From the update rule (1) and Lemma~\ref{lem:three} with $a=w_{t+1}$, $b=w_{t}$ and $c=w^{*}$,
\[
\|w_{t+1}-w^{*}\|^{2}
= \|w_{t}-w^{*}\|^{2}
-2\eta_{t}\langle Q(g_{t}),\,w_{t}-w^{*}\rangle
+\eta_{t}^{2}\|Q(g_{t})\|^{2}. \tag{3}
\]
\textbf{[Step 1]}

\subsection*{Step 2: Take Conditional Expectation}
Condition on $\mathcal{F}_{t}:=\sigma(w_{0},\xi_{0},\dots,\xi_{t-1})$. Using unbiasedness of $Q$ (Assumption~\ref{ass:quant}) we have
\[
\E\!\big[\langle Q(g_{t}),w_{t}-w^{*}\rangle\mid\mathcal{F}_{t}\big]
= \langle \E[Q(g_{t})\mid\mathcal{F}_{t}],\,w_{t}-w^{*}\rangle
= \langle g_{t},w_{t}-w^{*}\rangle.
\]
Similarly, by the variance bound,
\[
\E\!\big[\|Q(g_{t})\|^{2}\mid\mathcal{F}_{t}\big]
= \E\!\big[\|g_{t}+ (Q(g_{t})-g_{t})\|^{2}\mid\mathcal{F}_{t}\big]
= \|g_{t}\|^{2}+\E\!\big[\|Q(g_{t})-g_{t}\|^{2}\mid\mathcal{F}_{t}\big]
\le (1+\omega)\|g_{t}\|^{2}. \tag{4}
\]
\textbf{[Step 2]}

Taking $\E[\cdot\mid\mathcal{F}_{t}]$ on (3) and using (4) gives
\[
\E\!\big[\|w_{t+1}-w^{*}\|^{2}\mid\mathcal{F}_{t}\big]
\le \|w_{t}-w^{*}\|^{2}
-2\eta_{t}\langle g_{t},w_{t}-w^{*}\rangle
+\eta_{t}^{2}(1+\omega)\|g_{t}\|^{2}. \tag{5}
\]
\textbf{[Step 3]}

\subsection*{Step 3: Insert Stochastic Gradient Decomposition}
Write $g_{t}= \nabla f(w_{t})+\delta_{t}$ where $\delta_{t}:=g_{t}-\nabla f(w_{t})$ satisfies $\E[\delta_{t}\mid\mathcal{F}_{t}]=0$ and $\E[\|\delta_{t}\|^{2}\mid\mathcal{F}_{t}]\le\sigma^{2}$ (Assumption~\ref{ass:sgd}).  
Plugging into (5) and expanding the inner product:
\[
-2\eta_{t}\langle g_{t},w_{t}-w^{*}\rangle
= -2\eta_{t}\langle\nabla f(w_{t}),w_{t}-w^{*}\rangle
-2\eta_{t}\langle\delta_{t},w_{t}-w^{*}\rangle. \tag{6}
\]
Similarly,
\[
\|g_{t}\|^{2}= \|\nabla f(w_{t})\|^{2}+2\langle\nabla f(w_{t}),\delta_{t}\rangle+\|\delta_{t}\|^{2}. \tag{7}
\]

Taking full expectation and using $\E[\langle\delta_{t},w_{t}-w^{*}\rangle]=0$,
\[
\E\!\big[\|w_{t+1}-w^{*}\|^{2}\big]
\le \E\!\big[\|w_{t}-w^{*}\|^{2}\big]
-2\eta_{t}\E\!\big[\langle\nabla f(w_{t}),w_{t}-w^{*}\rangle\big]
+\eta_{t}^{2}(1+\omega)\big(\E\!\big[\|\nabla f(w_{t})\|^{2}\big]+\sigma^{2}\big). \tag{8}
\]
\textbf{[Step 4]}

\subsection*{Step 4: Use Convexity}
Convexity of $f$ gives
\[
f(w_{t})-f(w^{*})\le \langle\nabla f(w_{t}),w_{t}-w^{*}\rangle. \tag{9}
\]
Insert (9) into (8):
\[
\E\!\big[\|w_{t+1}-w^{*}\|^{2}\big]
\le \E\!\big[\|w_{t}-w^{*}\|^{2}\big]
-2\eta_{t}\E\!\big[f(w_{t})-f(w^{*})\big]
+\eta_{t}^{2}(1+\omega)\big(\E\!\big[\|\nabla f(w_{t})\|^{2}\big]+\sigma^{2}\big). \tag{10}
\]
\textbf{[Step 5]}

\subsection*{Step 5: Relate Gradient Norm to Function Suboptimality}
For $L$‑smooth convex functions,
\[
\|\nabla f(w_{t})\|^{2}\le 2L\big(f(w_{t})-f(w^{*})\big). \tag{11}
\]
(Proof: apply smoothness inequality with $y=w^{*}$ and rearrange.)  
Substituting (11) into (10) yields
\[
\E\!\big[\|w_{t+1}-w^{*}\|^{2}\big]
\le \E\!\big[\|w_{t}-w^{*}\|^{2}\big]
-2\eta_{t}\E\!\big[f(w_{t})-f(w^{*})\big]
+2L\eta_{t}^{2}(1+\omega)\E\!\big[f(w_{t})-f(w^{*})\big]
+ \eta_{t}^{2}(1+\omega)\sigma^{2}. \tag{12}
\]
\textbf{[Step 6]}

\subsection*{Step 6: Collect the Suboptimality Terms}
Define $\Delta_{t}:=\E[f(w_{t})-f(w^{*})]$.  
Equation (12) becomes
\[
\E\!\big[\|w_{t+1}-w^{*}\|^{2}\big]
\le \E\!\big[\|w_{t}-w^{*}\|^{2}\big]
-2\eta_{t}\Delta_{t}+2L\eta_{t}^{2}(1+\omega)\Delta_{t}
+\eta_{t}^{2}(1+\omega)\sigma^{2}. \tag{13}
\]
Rearrange:
\[
2\eta_{t}\Delta_{t}
\le \E\!\big[\|w_{t}-w^{*}\|^{2}\big]-\E\!\big[\|w_{t+1}-w^{*}\|^{2}\big]
+2L\eta_{t}^{2}(1+\omega)\Delta_{t}
+\eta_{t}^{2}(1+\omega)\sigma^{2}. \tag{14}
\]
\textbf{[Step 7]}

\subsection*{Step 7: Choose Stepsize Small Enough}
Assume $\eta_{t}\le \frac{1}{2L(1+\omega)}$, which holds for the schedule $\eta_{t}=c/\sqrt{t+1}$ with a sufficiently small constant $c$.  
Then $2L\eta_{t}^{2}(1+\omega)\le \eta_{t}$, and (14) simplifies to
\[
\eta_{t}\Delta_{t}
\le \frac12\Big(\E\!\big[\|w_{t}-w^{*}\|^{2}\big]-\E\!\big[\|w_{t+1}-w^{*}\|^{2}\big]\Big)
+\frac12\,\eta_{t}^{2}(1+\omega)\sigma^{2}. \tag{15}
\]
\textbf{[Step 8]}

\subsection*{Step 8: Sum Over Iterations}
Sum (15) for $t=0,\dots,T-1$:
\[
\sum_{t=0}^{T-1}\eta_{t}\Delta_{t}
\le \frac12\Big(\|w_{0}-w^{*}\|^{2}-\E\!\big[\|w_{T}-w^{*}\|^{2}\big]\Big)
+\frac12(1+\omega)\sigma^{2}\sum_{t=0}^{T-1}\eta_{t}^{2}. \tag{16}
\]
Discard the non‑negative term $-\E[\|w_{T}-w^{*}\|^{2}]$ and bound the sum of squares using $\eta_{t}=c/\sqrt{t+1}$:
\[
\sum_{t=0}^{T-1}\eta_{t}^{2}
= c^{2}\sum_{t=0}^{T-1}\frac{1}{t+1}
\le c^{2}\big(1+\log T\big)\le 2c^{2}\log T\quad\text{(for $T\ge2$)}. \tag{17}
\]
However, for the purpose of the $O(1/\sqrt{T})$ rate we use the looser bound $\sum_{t=0}^{T-1}\eta_{t}^{2}\le c^{2}\log T\le c^{2}\sqrt{T}$ (since $\log T\le\sqrt{T}$ for $T\ge 1$).  
Thus
\[
\sum_{t=0}^{T-1}\eta_{t}\Delta_{t}
\le \frac{\|w_{0}-w^{*}\|^{2}}{2}
+\frac12(1+\omega)\sigma^{2}c^{2}\sqrt{T}. \tag{18}
\]

\subsection*{Step 9: Relate to the Averaged Iterate}
Because $\eta_{t}=c/\sqrt{t+1}$ is decreasing,
\[
\frac{1}{T}\sum_{t=0}^{T-1}\Delta_{t}
\ge \frac{1}{\sum_{t=0}^{T-1}\eta_{t}}\sum_{t=0}^{T-1}\eta_{t}\Delta_{t}.
\]
Since $\sum_{t=0}^{T-1}\eta_{t}=c\sum_{t=0}^{T-1}(t+1)^{-1/2}\ge 2c(\sqrt{T}-1)$,
\[
\frac{1}{T}\sum_{t=0}^{T-1}\Delta_{t}
\le \frac{ \|w_{0}-w^{*}\|^{2}+ (1+\omega)\sigma^{2}c^{2}\sqrt{T}}{4c\sqrt{T}}.
\]
Finally, by convexity of $f$, $f(\bar w_{T})\le \frac1T\sum_{t=0}^{T-1}f(w_{t})$, and taking expectations yields
\[
\E\!\big[f(\bar w_{T})-f(w^{*})\big]
\le \frac{\|w_{0}-w^{*}\|^{2}}{2c\sqrt{T}}
+\frac{cL}{2\sqrt{T}}\big(\sigma^{2}+ \omega G^{2}\big),
\]
where we used $G^{2}\ge \sup_{t}\E[\|\nabla f(w_{t})\|^{2}]\le 2L\,\sup_{t}\Delta_{t}$ to replace $\|\nabla f(w_{t})\|^{2}$ by a uniform bound. This is exactly the bound (2). \qedhere

\textbf{[Step 9]}

%======================================================================
\section*{Rate Derivation (With Constants)}
%======================================================================
The bound (2) can be written as
\[
\E\!\big[f(\bar w_{T})-f(w^{*})\big]\le
\underbrace{\frac{\|w_{0}-w^{*}\|^{2}}{2c}}_{A}\frac{1}{\sqrt{T}}
\;+\;
\underbrace{\frac{cL}{2}\big(\sigma^{2}+ \omega G^{2}\big)}_{B}\frac{1}{\sqrt{T}}.
\]
Choosing $c=\frac{\|w_{0}-w^{*}\|}{\sqrt{L(\sigma^{2}+\omega G^{2})}}$ balances the two terms and yields
\[
\E\!\big[f(\bar w_{T})-f(w^{*})\big]\le
\frac{\|w_{0}-w^{*}\|\sqrt{L(\sigma^{2}+\omega G^{2})}}{\sqrt{T}}.
\]
Thus the optimal constant in front of $1/\sqrt{T}$ is proportional to $\sqrt{L(\sigma^{2}+\omega G^{2})}$.

%======================================================================
\section*{Special Cases}
%======================================================================
\begin{enumerate}
\item \textbf{Exact gradients, no quantization} ($\sigma=0$, $\omega=0$):  
The bound reduces to $\frac{L\|w_{0}-w^{*}\|^{2}}{2c\sqrt{T}}$, the classic deterministic gradient descent rate.

\item \textbf{No quantization, stochastic gradients} ($\omega=0$):  
We recover the standard SGD bound $\mathcal{O}\!\big((\sigma^{2}L)^{1/2} / \sqrt{T}\big)$.

\item \textbf{Deterministic gradients, quantized communication} ($\sigma=0$):  
The error is driven solely by the quantization variance:
\[
\E\!\big[f(\bar w_{T})-f(w^{*})\big]\le
\frac{\|w_{0}-w^{*}\|^{2}}{2c\sqrt{T}}+
\frac{cL\omega G^{2}}{2\sqrt{T}}.
\]
\end{enumerate}

\end{document}